\subsection{Coherent backgrounds and the law topology}
\label{tree:sec:coherent-background}

We next put all scales in one space of observations.  Its purpose is
to retain the same continuum background under time changes and to
make the later success tests stable under extraction.

\paragraph{Retained observations.}
Enumerate the
bounded rational time windows, the finite Poisson/rational chart
instructions of Section~\ref{tree:sec:certificates}, their rational
open margins and cost cutoffs, and all finite lists of subpath
marks.  For each instruction use its compact closed-certificate
hyperspace from Proposition~\ref{tree:prop:certificate_limit}.  Record
also the internal endpoint variables in $[0,\infty]$, the
continuous contour on compact windows, and the continuum
background used in Proposition~\ref{tree:prop:zoom_transport}.
The background includes the curve and reference metric in local
uniform topology and the fields in their distribution topologies.  It is
given its usual standard Polish realization.  For every bounded
connected rational coordinate domain $W$ and every compact
rational exhaustion set $K\Subset W$, retain also
$D_h(\cdot,\cdot;W)|_{K\times K}$ in uniform topology.
These internal-reference-metric observations have a fixed
joint marginal; retaining them does not assert continuity
of the map from a field or an ambient metric to its internal
metric.  We also retain,
as deterministic coordinates, the joint laws of every auxiliary
finite reference experiment in
Lemma~\ref{tree:lem:finite_kernel_transfer}: the canonical input
and complete marked output law for each package, all finite
lists of independent packages, and their common refinements.
Use the same prescribed smooth input densities at every scale.
Include the countably many fixed quantum-unit transforms of
these templates.  At finite resolution the input marginals
depend on the random-walk resolution and offset lattice;
the invariance principle and the lattice-cell approximation of
forest offsets make them converge to the prescribed continuum input
marginals.  These limiting
input marginals are the same for every continuum law in a
subsequent varying-law extraction.  Each output space is a
compact hyperspace, so input tightness gives precompactness
of the finite-resolution weak-law coordinates as well.
They are auxiliary reference
laws, not extra samples from the actual contour.  Their
compatibility with the actual atlas, including the
density-and-bridge identities, is retained by
Lemma~\ref{tree:lem:certificate_class_closed}.

A bounded product
metric, with coordinate weights $2^{-r}$, makes this a separable
metric data space $\mathcal E$; use the closed compatible-data
subspace when a refinement identity is imposed.  We equip its
laws with a metric $\mathbf d$ for weak convergence, for example
the bounded Lipschitz metric for this bounded product metric.
Equivalently, this topology is specified by integrals against a
fixed countable convergence-determining family of bounded
continuous functions.  Thus all tests below refer to one
specified countable law topology.

The raw certificate coordinates are the time-labelled compact
spaces in Equation~\eqref{tree:eq:marked_certificate_space}.  Discrete
grid times are points of the same compact time window, so a
finite-resolution law belongs to this ambient space without
satisfying the Brownian quotient identities.  All choices of
representatives are retained.  Spatial projection is applied
only after adjoining the continuum background; the vanishing
projected diameter of individual links supplies its limiting
compatibility.

\paragraph{One Brownian coupling for all scales.}
Take independent
two-sided Brownian motions $W_1,W_2$ and an independent sequence
of fair coordinate choices.  For each coordinate, read successive
simple-walk increments from the signs of Brownian exits from the
interval of radius one about its current position.  Use the
coordinate choices to interleave the two resulting walks; apply
the same construction independently on the negative half-line.
This gives exactly the bilateral axial walk.  The successive
Brownian exit durations are independent, have mean one, and
satisfy the law of large numbers.  The numbers of each coordinate
choice likewise have asymptotic frequency $1/2$.  Hence the
Brownian clock used by coordinate $a$ at walk time $Nt$, divided
by $N$, converges uniformly on bounded $t$-intervals to $t/2$,
in probability.  The Brownian modulus of continuity after
diffusive rescaling implies
\begin{equation}
 \begin{gathered}
 \sup_{|t|\leq T}N^{-1/2}
 \bigl|Z^{\rm disc}(\lfloor Nt\rfloor)-B(Nt)\bigr|
 \longrightarrow0\quad\text{in probability},\\
 B(t)=(W_1(t/2),W_2(t/2)),
 \end{gathered}
 \label{tree:eq:bootstrap_common_coupling}
\end{equation}
for every fixed $T<\infty$.  The interpolation error is at most
one step.  This argument only uses the strong Markov property,
the exit-time mean, and the law of large numbers.

\paragraph{Attach and transport the continuum background.}
Let $\nu$ be the law of the full bilateral Brownian contour.
Write $\nu(db)\mathsf Q(b,dc)$ for its standard joint law with
the full continuum background $\mathcal C$, including any
residual imaginary-field and embedding data.  This background
excludes the microscopic coordinate schedule, time cuts, and
additional passage data.  If
$S_vb(t)=v^{-1/2}b(vt)$, its intrinsic quantum-unit transform is
\[
 T_v\mathcal C=(h-c_v,\widehat h,\eta(v\,\cdot),v^{-p}D_h),
 \qquad c_v=\gamma^{-1}\log v,
\]
followed by the prescribed common reembedding of its components.
The decorated-cone law is invariant under $(S_v,T_v)$ by
Equation~\eqref{tree:eq:background_joint_zoom}.
Here the imaginary field has the uniform phase modulo its
angular period specified in
Proposition~\ref{tree:prop:background_locality}.  Under a positive
circle-average dilation, its GFF modulo constants is invariant,
and the change of anchoring constant is absorbed by that
uniform phase.  The dilation is determined by the real field,
which is independent of the imaginary field before the curve
is drawn.  The curve then transforms by imaginary-geometry
covariance, and the reference metric by its established
covariance and constant-shift rule.  Thus the invariance is
for the full background, including the extra imaginary
randomness.  A measurable representative is chosen from the
intrinsic decorated surface and never from passage data.

Disintegration of this joint invariance gives
\begin{equation}
 \mathsf Q(S_vb,\cdot)=(T_v)_\#\mathsf Q(b,\cdot)
 \quad\text{for $\nu$-almost every $b$}.
 \label{tree:eq:bootstrap_background_kernel_equivariance}
\end{equation}
For the countable dyadic group, take one full-measure set on
which these identities and their compositions all hold.  Draw
$\mathcal C$ once from $\mathsf Q(B,\cdot)$, independently of
the coordinate schedule and all time-cut processes conditionally
on $B$.  Put $B_j=S_{2^j}B$ and $\mathcal C_j=T_{2^j}\mathcal C$.
The background at $j-i$ is exactly the time-$2^{-i}$ transport
of the background at $j$, intrinsically.  The discrete
instruction data obey the same identity, with the extra
deterministic cost conversion $A_{j-i}/A_j$.

Let $Y_j$ denote the raw, unprojected time-labelled passage
data and time cuts at scale $j$.  Since $S_{2^j}$ is invertible
on the full bilateral contour, $B_j$ determines $B$.
The conditional sampling and
\eqref{tree:eq:bootstrap_background_kernel_equivariance} give the
exact finite-resolution identity
\begin{equation}
 \Law(\mathcal C_j\mid B_j,Y_j)=\mathsf Q(B_j,\cdot).
 \label{tree:eq:bootstrap_background_conditional}
\end{equation}
One must use raw data here: their projected spatial ranges
also depend on $\mathcal C_j$ and are adjoined afterward.

\paragraph{Preserve the conditional law under extraction.}
We verify that \eqref{tree:eq:bootstrap_background_conditional}
passes through extraction despite possible discontinuity of
$\mathsf Q$.  Suppose $(B_j,Y_j)\Rightarrow(B_\infty,Y_\infty)$.
The $B_j$ marginal is always $\nu$, and the $\mathcal C_j$
marginal is always $\nu\mathsf Q$, so the joint triples are
tight.  For bounded continuous $f$ and $g$, write
$\mathsf Qg(b)=\int g(c)\mathsf Q(b,dc)$.  Then
\[
 \E[f(B_j,Y_j)g(\mathcal C_j)]
 =\E[f(B_j,Y_j)\mathsf Qg(B_j)].
\]
Approximate $\mathsf Qg$ in $L^1(\nu)$ by bounded continuous
functions with a common bound.  The replacement error is at
most $\|f\|_\infty$ times this $L^1$ error, uniformly in $j$
and also in the limit.  Weak convergence applies to each
continuous approximation.  Removing the approximation yields
\begin{equation}
 \Law(\mathcal C_\infty\mid B_\infty,Y_\infty)
       =\mathsf Q(B_\infty,\cdot).
 \label{tree:eq:bootstrap_background_closure}
\end{equation}
Products of continuous tests determine the joint law, giving
the stated conditional identity.  Thus conditioning the
limiting reference experiments additionally on the continuum
fields does not change their passage kernels.  This is a
fixed-kernel argument; closure when the passage kernels vary
uses the retained reference-law coordinates and
Lemma~\ref{tree:lem:certificate_class_closed} instead.

\paragraph{The law at scale $j$.}
Write $\mu_j$ for the law of these complete data in time units
$2^j$ and cost units $A_j$.  Equation~\eqref{tree:eq:bootstrap_common_coupling}
and Proposition~\ref{tree:prop:certificate_limit} give the required
joint extractions.  Recording the continuum background avoids
asserting that an arbitrary measurable functional of a contour
is continuous.  Its joint marginal is fixed, and its contour
equals the discrete contour limit.  The exact transport
identities just established identify all fixed-offset limits
with the corresponding zooms in a common extraction.

Fixed local changes of density can be included in this topology
without a new compactness assumption.  Truncate each density and
approximate it in $L^1$ of the fixed background marginal by bounded
continuous functions.  Weak convergence applies to the
approximations, and their $L^1$ errors are uniform over the
coupled laws.  Use a countable list of the buffered domains and
approximations required for the shell tests.  In particular, the
local free-field tests below are tests of the same kernels
identified in Proposition~\ref{tree:prop:locality}.

